\documentclass{article}
\usepackage[T1]{fontenc}
\usepackage{lmodern}
\usepackage{graphicx}
\usepackage{amsmath,amssymb}
\usepackage[a4paper,margin=2cm]{geometry}
\usepackage[absolute,overlay]{textpos}
\setlength{\TPHorizModule}{1mm}
\setlength{\TPVertModule}{1mm}
\usepackage[indonesian]{babel}
\usepackage{hyperref}
\setlength{\emergencystretch}{2em}
% unit: Halaman 107

\begin{document}

% === Halaman 107 (python/native) ===

107

• Hill cipher mudah dipecahkan dengan known-plaintext attack.

• Misalkan untuk Hill cipher dengan m = 2 diketahui:

• P = (19, 7)    → C = (0, 23)

• P = (4, 17)    → C = (12, 6)

• Jadi, K(19, 7) = (0, 23) dan K(4, 17) = (12, 6)

• Matriks balikan dari P adalah P-1 =  





• Sehingga




→



→















=













P 


C

26


mod


17

7

4

19

K


6

23

12

0













=













−

5

5

14

25



26


mod


17

7

4

19

1













=

























14

7

8

8



26


mod


5

5

14

25


6

23

12

0

Kriptanalisis Hill Cipher

K = CP–1 mod 26 = 

→ K = CP–1 mod 26  

Enkripsi: C = KP mod 26

Dekripsi: P = K-1C  mod 26

\end{document}
